\chapter{Head Lice}

\section*{Definition and Overview}
Head lice are tiny wingless insects that live on the scalp and feed on blood. They are very common, especially in primary school children, and spread through direct head-to-head contact. Head lice do not jump or fly, and they do not transmit disease, but they cause itching and can be a source of embarrassment and disruption at school. The main sign of infestation is itching of the scalp, and the diagnosis is confirmed by finding live lice or eggs (nits) attached to the hair. Head lice are treated with specially formulated lotions or with thorough wet combing, and they are entirely preventable and manageable.

\section*{Epidemiology}
Head lice are one of the most common childhood conditions, affecting an estimated 1 in 10 primary school children at any time. They are most common in children aged 4--12, and girls are affected more often than boys, partly because of longer hair. Head lice spread easily wherever children are in close contact, including schools, childcare, and sleepovers. Infestations are not related to poor hygiene; lice prefer clean hair. Most children have head lice at some stage, and repeat infestations are common because eggs can hatch and contacts are not always treated.

\section*{Aetiology and Pathophysiology}
Head lice are insects that live their whole life on the human scalp.

\begin{itemize}
    \item \textbf{The louse:} \textit{Pediculus humanus capitis} is a grey-brown insect about the size of a sesame seed
    \item \textbf{Transmission:} by direct head-to-head contact; lice crawl from one head to another. They do not jump, fly, or survive long away from the scalp
    \item \textbf{Life cycle:} the female lays eggs (nits) glued to the hair near the scalp; eggs hatch in about 7--10 days, and nymphs mature into adults in around 10 days
    \item \textbf{Feeding:} lice feed on blood from the scalp several times a day
    \item \textbf{Itching:} an allergic reaction to louse saliva causes itching, which can take weeks to develop
    \item \textbf{Survival:} lice need a human host; they cannot live on pets or furniture
\end{itemize}

\section*{Clinical Features}
The features of head lice infestation include:

\begin{itemize}
    \item Itching of the scalp, especially behind the ears and at the back of the neck
    \item Live lice, which move quickly and can be hard to see
    \item Nits: small, oval, grey-white eggs glued to the hair shaft near the scalp
    \item A scratchy feeling or irritation
    \item Sores and crusting on the scalp from scratching
    \item Swollen lymph nodes behind the ears in some cases
    \item Many children have no symptoms at all
    \item Empty nit shells may remain after treatment
\end{itemize}

\section*{Differential Diagnosis}
Other scalp conditions can be mistaken for head lice.

\begin{itemize}
    \item Dandruff and seborrhoeic dermatitis, which cause flakes that are easily brushed off (unlike nits, which are firmly attached)
    \item Dry hair and hair product residue
    \item Scabies, which causes itching and a rash elsewhere on the body
    \item Flea bites and insect bites
    \item Eczema of the scalp
    \item Hair casts and other hair shaft conditions
    \item A fear of lice when there are no lice (imagined infestation)
\end{itemize}

\section*{When to Seek Medical Care}
Head lice do not usually need medical care and are treated at home. See a pharmacist for advice on treatments, and a doctor if the scalp becomes infected, if the infestation persists despite treatment, or for advice for babies, pregnant women, or people with sensitive skin.

\textbf{Call 000 (or local emergency services) for:}
\begin{itemize}
    \item This condition is not an emergency; however, seek care for spreading redness, swelling, pain, or fever from a secondary scalp infection
    \item Any severe allergic reaction to a head lice treatment
\end{itemize}

\section*{Diagnosis and Investigations}
Diagnosis is made by finding live lice or nits.

\begin{itemize}
    \item \textbf{Wet combing:} combing the hair with a fine-toothed nit comb after applying conditioner makes live lice visible
    \item \textbf{Examination of nits:} nits are glued to the hair, unlike dandruff, and are found close to the scalp in active infestation
    \item \textbf{Distinguishing live from dead nits:} live nits are dark and close to the scalp; empty shells are white and further out
    \item \textbf{Checking household contacts:} family members are checked and treated if affected
    \item \textbf{Follow-up:} repeat checks confirm the lice are gone
\end{itemize}

\section*{Management}
Treatment is effective when applied correctly and combined with checking contacts.

\begin{itemize}
    \item \textbf{Permethrin or other insecticide lotions:} applied to dry hair, left for the recommended time, then washed out; a second application is usually needed after 7--10 days
    \item \textbf{Dimeticone lotion:} a silicone-based product that suffocates lice, applied to dry hair and left overnight or for the recommended time
    \item \textbf{Wet combing:} a thorough method using conditioner and a fine-toothed comb every 2--3 days for at least two weeks
    \item \textbf{Combing out nits:} removing nits with the comb after treatment
    \item \textbf{Treating contacts:} checking and treating family members with live lice
    \item \textbf{Notifying the school:} informing the school so other families can check
    \item \textbf{Follow-up checks:} confirming the infestation has cleared
    \item \textbf{No need to fumigate:} lice do not survive long off the head, so extensive cleaning is not required
\end{itemize}

\section*{Complications}
\begin{itemize}
    \item Secondary bacterial infection of scratched skin, requiring treatment
    \item Persistent infestation when treatment is not applied correctly or contacts are not treated
    \item Repeated infestations, which are common
    \item Skin irritation from treatments
    \item Anxiety, stigma, and missed school days
    \item Misuse of treatments or home remedies, which can be ineffective or harmful
\end{itemize}

\section*{Prognosis and Outcomes}
Head lice are completely treatable. With correct treatment and thorough combing, the infestation clears within one to two weeks, and recurrence is prevented by checking and treating contacts. Head lice do not cause serious illness or transmit disease. Most families manage infestations successfully at home, and children can return to school once treatment has been completed. Although lice are a nuisance, the outlook is excellent.

\section*{Prevention and Self-Care}
\begin{itemize}
    \item Check children's hair regularly, especially after known exposure
    \item Treat promptly when live lice are found
    \item Follow the treatment instructions exactly, including the second application
    \item Comb with a fine-toothed nit comb during and after treatment
    \item Check and treat all household members with live lice at the same time
    \item Do not share combs, brushes, hats, or hair accessories
    \item Do not use harsh home remedies or multiple treatments at once
    \item Inform the school or childcare centre so other families can check
    \item Do not keep children home unnecessarily; they can return after treatment
\end{itemize}

\section*{Sources and Further Reading}
The following reputable sources provide further information for healthcare students and patients seeking reliable, current advice about head lice.

\begin{itemize}
    \item Healthdirect Australia. \textit{Head lice.} https://www.healthdirect.gov.au/
    \item Australian Government Department of Health and Aged Care. \textit{Head lice fact sheet.} https://www.health.gov.au/
    \item NSW Health. \textit{Head lice fact sheet.} https://www.health.nsw.gov.au/
    \item Raising Children Network (Australian Government). \textit{Head lice.} https://raisingchildren.net.au/
    \item NHS. \textit{Head lice and nits.} https://www.nhs.uk/
    \item Centers for Disease Control and Prevention. \textit{Head lice.} https://www.cdc.gov/
\end{itemize}
